% \section{Contribution II: An entity-centric approach to manage court judgments based on Natural Language Processing - Entity Linking for Italian Legal Judgments}
\section[Entity Extraction from Italian Civil Judgments][EE from Italian Judgments]{Entity Extraction from Italian Civil Judgments}
% \sectionmark{EE from Italian Judgments}
\label{sec:eejudgements}
% \todo{maybe integrate commented code here to legal sota or preface}

In this section, we evaluate the performance of existing \acf{ee} systems on Italian civil court judgments.
The work presented in this section is related to the publication \textcite{BELLANDI2024}.

As discussed in \Cref{sec:background:ai4legal}, legal data often contain personal information, making compliance with privacy regulations a key challenge. This is reflected in the scarcity of labeled legal data for training or evaluating \ac{ee} systems~\cite{li-etal-2022-parameter}. Furthermore, the language used in legal documents differs significantly from everyday language~\cite{ccetindaug2023name_legal_turkish}, whereas most open datasets employed to train neural models---such as Common Crawl~\cite{commoncrawl} or The Pile~\cite{gao2020pile800gbdatasetdiverse}---are derived from general web sources.

Another related issue concerns the use of external \acp{api}, such as those offered by OpenAI\footnote{\url{https://openai.com/api/}}, on legal data containing personal information. Unless the input data are anonymized or pseudonymized, their use may result in transferring sensitive data outside the European Union, raising compliance concerns with respect to the GDPR regulations~\cite[Art.~44--49]{gdpr2016}. Consequently, many institutions must rely on locally deployed models, which, while ensuring full control over data storage and security~\cite[Art.~32]{gdpr2016}, may be constrained by the available computational infrastructure.

In summary, the main contributions of this work on \acl{ee} applied to Italian civil judgments are as follows:
\begin{itemize}
    \item Creation of an annotated benchmark dataset for civil judgments, including the calculation of the \acf{iaa}, for enabling structured evaluation of \acl{ee} methods.
    \item Evaluation of \acl{ner} and \acf{inel} pipelines on domain-specific data (Italian civil judgments), assessing performance and identifying key challenges.
    \item Comparative study of domain adaptation techniques, highlighting which approaches provide the best trade-off between computational cost and performance improvement.
    \item Training of the BLINK~\cite{blink} bi-encoder for the Italian language (\blinkita).
\end{itemize}

The remainder of this section is organized as follows: \Cref{sec:judgements_labeled_data} describes the creation of the labeled benchmark dataset; \Cref{sec:italia_ee_pipeline} presents the Italian \ac{ee} pipeline; and \Cref{sec:judgment_ee_eval} reports the experimental evaluation.

\subsection{Annotation of a Legal Benchmark}
\label{sec:judgements_labeled_data}

We started from a corpus of 927,453 judgments in civil trials published from 2008 to 2021 (the majority of them, $\approx 86\%$, have been published from 2016 to 2021).
This number can be compared to the total number of trials that are estimated per year according to~\cite{bankitalia}, which falls between 2 and 2.5 millions from 2010 to 2019.
For the remainder of this work, we will refer to this corpus or more than 900 thousand civil judgments as \emph{ICCJ900k}---from ``Italian Civil Court Judgments''.

Data includes the judgment text and 41 metadata with information about the judge (or the president if several judges are involved), the number and year of the judgment and of the trial, the court and the district it belongs to, the instance (trial or appeal), references to the trial in case of appeal, a code describing the subject and additional technical fields. This dataset has been provided by the Italian Ministry of Justice and consists of real documents as they were archived. As a consequence, in many cases, the texts contain spurious lines that must be cleared before processing, for instance, some metadata at the very beginning, duplicate headings, extra blank lines or characters from stamps present in the printed version.

%structure
The judgments' structure consists of several sections. The most important are the following:
\begin{itemize}
    \item the preamble with the judge(s), plaintiff(s) and defendant(s) data,
    \item the description of the case,
    \item and the final decision and dispositions with the related reasons.
\end{itemize}

%%%%%%%

From ICCJ900k, we constructed a benchmark dataset of 146 documents labeled for \acf{ee}---which we call \emph{ICCJ146-EE}. The documents have been selected using stratified sampling on the province of the court that made the judgment.
The annotation steps was performed semi-automatically according to the following order:
\begin{enumerate}
    \item the documents have been automatically annotated for \ac{ner};
    \item human annotators corrected \ac{ner} errors;
    \item \ac{nel} with NIL prediction methods have been executed on corrected \ac{ner} annotations;
    \item human annotators corrected \ac{nel} errors. For NIL mentions, they also assigned a novel entity name to each NIL mention so that all mentions referring to the same new entity could be grouped under the same label, obtaining also NIL clustering annotations.
\end{enumerate}
In summary, the annotations include labels for \ac{ner}, \ac{nel}, NIL prediction, and NIL clustering (see \Cref{sec:incrementalel} for details on these tasks).

This approach, where algorithms are used to speed up annotation has been used for labeling \ac{ner} data in previous work~\cite{lrec2022kind}, and we can consider this process as a \say{simulation} of a real \emph{\acf{hitl}} use case, where humans manually verify and correct automatic annotations, for enhancing their quality.
Furthermore, this \say{simulation} allows to estimate the amount of time required by the final users-in-the-loop to correct automatic annotations obtained with the \ac{ee} pipeline described in \Cref{sec:italia_ee_pipeline}.

Initially, \ac{ner} annotations have been labeled on all the 146 documents by two annotators. Later, a subset of 30 documents has been additionally labeled for \ac{inel}. We refer to this subset as ``ICCJ30-\ac{inel}''.
The annotation process focused the entities of class \texttt{Person}, \texttt{Location}, \texttt{Organization}, \texttt{Money}, \texttt{Date}, and \texttt{Miscellaneous}.


To obtain an objective measure of the quality of the \ac{ner} annotations, which served as the starting point for subsequent annotation processes, we calculated an \acf{iaa} measure.

\paragraph{\ac{ner} Annotation Process}
An \ac{ner} label---as described in \Cref{sec:rw:ner}---is composed of the indexes of start and end characters, and by the type of the mentioned entity $(i_{start},i_{end},t^i)$.
The annotations were carried out by two annotators using \textit{doccano}~\cite{doccano} a web-based application for labeling tasks. The annotator was initially assigned with two disjoint sets of 73 documents. Then, from each of the two set we selected 15 documents and assigned them to the other annotator, so that these $15+15$ documents were annotated by both the annotators. In the end, each annotator received 88 documents and the double-annotated 30 were used to measure the \ac{iaa}.  


Based on~\cite{jha2017all}, a detailed set of annotation guidelines was prepared and provided to the annotators for reducing the inconsistencies arising from different annotation styles. During the annotation process, the guidelines were iteratively refined by recording and addressing the annotators' doubts and edge cases.

The guidelines cover both span selection and type assignment.  
For span selection, for example in the case of nested entities, the span corresponding to the most informative and specific entity type should be annotated. ``\underline{Main Street 18, London}[LOC]'' should be preferred over ``\underline{Main Street 18}[LOC], \underline{London}[LOC]''.  

For type assignment, the entity type should be determined from the surrounding context. For instance, ``John Smith, born in \underline{France}[LOC]'' versus ``\underline{France}[ORG] announced new trade regulations''.  

Each entity type also has its own specific guidelines to handle particular exceptions. For example, titles should be excluded from person names: ``Dr. \underline{John Smith}[PER]'' is preferred over ``\underline{Dr. John Smith}[PER]''.

Finally, we also decided to annotate mentions such as ``the Judge'' or ``the Court'', which are not strictly within the scope of \ac{ner} but are closely related to coreference resolution, as these entities are particularly relevant in the legal domain.




\paragraph{Document Statistics}
The resulting corpus counts 16,634 annotations ($\approx$114 annotations per document), with an average document length of $\approx$1,900 words. Each document required on average $\approx$15 minutes to be annotated. The distribution of types is reported in \Cref{fig:types_distrib}.

\begin{figure}
    \centering
    \includegraphics[scale=0.8]{images/ICCJ146-EE_entity_type_distribution.png}
    \caption{Distribution of \ac{ner} types in the ICCJ146-EE dataset (with NIL counts).
    % *\textit{location} implicitly includes \textit{addresses} since nested annotations are not taken into account in this work; **miscellaneous annotations are mostly referred to \textit{law articles}, thus there is no distinction in the count.
    }
    \label{fig:types_distrib}
\end{figure}


\paragraph{Inter-Annotator Agreement}
To assess the quality of the annotations and the clarity of the guidelines we compute the \acl{iaa} on the 30 overlapping documents, which are annotated independently by both the annotators.
% These documents were sampled at different timestamps in order to be representative since the annotators became more confident over time.


Inter-annotator agreement measures such as \textit{Scott's~$\pi$}, \textit{Cohen's~$\kappa$}, \textit{Fleiss'~$\kappa$}, and \textit{Krippendorff's~$\alpha$} were developed to correct observed agreement for the portion that may occur purely by chance~\cite{vanOest2019}. In other words, these measures estimate how much of the observed consistency between annotators exceeds what would be expected if both raters assigned categories independently according to their own labeling probabilities~\cite{HayesKrippendorff2007}.
Formally, they compare the observed agreement $P_o$ to the expected agreement $P_e$ under a chance model, and express reliability with the following formula:
\begin{equation}
    (P_o - P_e)/(1 - P_e).
\end{equation}

However, these measures may not be suitable for complex labeling tasks such as \ac{ner}, providing low interpretability values~\cite{braylan2022measuring}. A problem of these measures for \ac{ner} is that we deal with text spans, so it is difficult to define a general criterion to identify positive and negative examples to calculate a consistent IAA value~\cite{deleger2012building}. A simple solution to face this problem is to compute the metrics at token-level, however this would yield overly optimistic IAA value. For example, ``\underline{Barack Obama}[PER] was born in Honolulu'' and ``\underline{Barack}[PER] \underline{Obama}[PER] was born in Honolulu'' would be considered as a perfect match and the high number of negatives---\ie the tokens outside of entity mentions---causes the calculation of the metrics on a very imbalanced data~\cite{deleger2012building}.

Also, the aforementioned IAA coefficients differ primarily in how the expected-by-chance term $P_e$ is estimated~\cite{vanOest2019}, with Fleiss'~$\kappa$ and Scott's~$\pi$ being equivalent when considering two annotators, since the former is a generalization of the latter~\cite{HayesKrippendorff2007}. Thus, when there are only two annotators and the category proportions are highly unequal---as is common in \ac{ner}, where most tokens belong to the ``outside-of-entity-mention'' class---both annotators exhibit almost identical marginal label distributions. Consequently, these coefficients tend to converge toward similar values in such settings.

For these reasons, we adopted the \acl{ner} F$_1$ score as the main metric, since it is a more robust alternative to the classic \ac{iaa} metrics~\cite{hripcsak2005agreement,grouin2011proposal}.
%, wang2021nero}.
The F$_1$ score---calculated with the matching criteria described in \Cref{sec:metrics}---is pair-wisely computed between annotators: for each combination, the labels from an annotator are used as the ground truth to evaluate the labels from another annotator; an average is computed in case of more than two annotators. 

\Cref{tab:iaa} shows all metrics averaged over documents. For completeness, we also report the token-level F$_1$ score for the \textit{strong-typed} match criteria and Scott's~$\pi$, Cohen's~$\kappa$, Fleiss'~$\kappa$, and Krippendorff's~$\alpha$.
We can see that \textit{strong-typed (instance-level)}, \ie the strictest metric, reaches 80.8\%, which can be considered satisfactory since it measures perfect matches between annotators. Moreover, by looking at the additional F$_1$ score-based metrics we can see that the values increase as we relax the matching criterion (partial is the least stringent). This proves that even in cases where the labels from the two annotators do not perfectly match, there is still a significant degree of overlap.
These differences between F$_1$ scores---calculated with harder and softer matching criteria---as well as the difference between \textit{instance-level} and \textit{token-level} \textit{strong-typed} measures, also point out that a perfect span selection was the most difficult part for the annotators.

\begin{table}[]
\centering
\begin{tabular}{l|l}
\toprule
\textbf{Metric}                     & \textbf{Value} \\\toprule
\multicolumn{2}{c}{F$_1$ score-based metrics} \\\midrule
Strong                              & 83.9           \\
Strong-typed (instance-level)       & 80.8           \\
Strong-typed (token-level)          & 88.6           \\
Approximate                         & 96.6           \\
Approximate-typed                   & 90.4           \\
Partial                             & 96.9           \\
Partial-typed                       & 90.6           \\\midrule
\multicolumn{2}{c}{Standard IAA metrics} \\\midrule
$\text{Scott's~}\pi \approx \text{Cohen's~}\kappa \approx \text{Fleiss'~}\kappa \approx \text{Krippendorff's~}\alpha$
 & 66.2           \\\bottomrule
\end{tabular}
\caption{Inter-annotator agreement for \ac{ner} computed on a subset of 30 documents from ICCJ146-EE. The description of each F$_1$ score-based metric can be found in \Cref{sec:metrics}. *based on strong-typed match.}
\label{tab:iaa}
\end{table}


\paragraph{Knowledge Consolidation}
The annotation for the knowledge consolidation process---\ie \ac{nel}, NIL prediction and clustering---has been performed semi-automatically from the \ac{ner} labels of the 30 documents used to compute the IAA. The guidelines, based on the assumptions in~\cite{jha2017all}, are simpler in this case and be summed up with the following two steps.

\begin{enumerate}
    \item If the mention refers to a known entity in the \ac{kr} (Italian Wikipedia and manually added entities) label it with the entity URI and mark the mention as $\neg$NIL.
    \item Otherwise, mark the mention as NIL, create a new entity, and label the mention with the new entity URI. This latter step is necessary to trace which NIL mentions refer to the same unknown entity (NIL clustering). 
\end{enumerate}

% In addition, other entity-level guidelines (\eg \textit{``Italia''} must be linked to \textit{``\url{https://it.wikipedia.org/wiki/Italia}''} and not to \textit{``\url{https://it.wikipedia.org/wiki/Italia_(regione_geografica)}''}) and type-level guidelines (\eg a NIL entity of type \textit{person} must be labeled \textit{``Name Surname''}, thus linking mentions like \textit{``Mario Rossi''} and \textit{``Rossi Mario''} to the same out of \ac{kr} entity) have been defined during the annotation process when particular cases were found. 

To enhance the efficiency of the annotation procedure an application was developed, incorporating both Wikipedia search API and a fuzzy search mechanism. On average documents required $\approx$15 minutes to be annotated for \ac{nel}, NIL prediction, and NIL clustering. Out of the 3,006 annotated mentions, 467 refer to Wikipedia entities, and 1,753 to NIL entities. The remaining 786 mentions, of class \texttt{Date} and \texttt{Money}, were not assigned to any entity.

% \paragraph{NIL clustering}

The 2,200 mentions linked to an entity are organized into a total of 1,025 clusters, with 211 clusters corresponding to Wikipedia entities.
The larger portion of 814 refers to NIL entities, as expected within the considered domain.

\subsection{Italian Incremental Entity Linking Pipeline}
\label{sec:italia_ee_pipeline}

Inspired by the \ac{inel} pipeline described in \Cref{sec:ijckg:pipeline} and by recent work on \ac{inel}~\cite{Heist2023nastylinker,kassner-etal-2022-edin} we develop a pipeline-based system that, first, orchestrates \ac{ner} and then the knowledge consolidation processes, namely \ac{nel}, NIL prediction, and NIL clustering.
For \ac{ner} we combine different methods to extract all the types annotated in ICCJ146-EE.

\Cref{tab:algotypes} illustrates the types supported by each \ac{ner} extractor, as well as by the \ac{nel}, NIL prediction, and NIL clustering algorithms. Notably, for the latter three tasks, we exclusively consider the types \texttt{Person}, \texttt{Location}, and \texttt{Organization}. This decision is based on the fact that \texttt{Date} and \texttt{Money} are not directly associated with any linkable entity. With respect to the \texttt{Miscellaneous} class, during an exploratory phase, we found that the quality of the prediction for this class---significantly lower than for the other classes---was not promising enough to proceed with the knowledge consolidation process. We further discuss this issue in the experimental evaluation.


% Our primary objective is to assess the performance of existing algorithms, trained on publicly available training data, on civil judgments.
% and to examine the performance gap that arises when we employ algorithms trained on ``conventional'' training data (typically news articles) for the more challenging data scenario (domain-specific language; text extracted from PDFs).

\begin{table}[H]
    \centering
    % \scalebox{.7}{
        \begin{tabular}{l|cccc|c|c|c}
            \multirow{2}{*}{Types}      & \multicolumn{4}{c|}{\textbf{\Acl{ner}}}              & \multirow{2}{*}{\textbf{\ac{nel}}}  & \multirow{2}{*}{\textbf{NIL pred.}}    & \multirow{2}{*}{\textbf{NIL clust.}} \\
                                    & \Spacy & Tint  & Trie\ac{ner}*  & Combination    &                       &                               & \\
            \midrule
            \texttt{Person}         & \ding{51} & \ding{51} & \ding{51}*    & \ding{51}   & \ding{51}        & \ding{51} & \ding{51}   \\
            \texttt{Location}       & \ding{51} & \ding{51} & --             & \ding{51}   & \ding{51}        & \ding{51} & \ding{51}   \\
            % + Address       & RAE    & -         & -         & -             & -           & Map         & Implicit   & Implicit \\
            \texttt{Organization}   & \ding{51} & \ding{51} & \ding{51}*    & \ding{51}   & \ding{51}        & \ding{51} & \ding{51}   \\
            \texttt{Money}          & --         & \ding{51} & --             & \ding{51}   &     --        & --          & --  \\
            \texttt{Date}           & --         & \ding{51} & --             & \ding{51}   &     --       & --          & --  \\
            \texttt{Miscellaneous}  & \ding{51} & --         & --             & \ding{51}   &     --        & --          & --  \\
            % + Law article   & RAE    & -         & -         & -             & -           & Normattiva  & Implicit          & Implicit  \\
        \end{tabular}
        % }
    \caption{Types supported by the \ac{ner}, \ac{nel}, NIL prediction, and NIL Clustering algorithms. *Note that the Trie\ac{ner} relies on a gazetteer and it is potentially capable of extracting any entity regardless of its type as long as a pattern is provided. In our experiments, Trie\ac{ner} extracts entities of type \texttt{Person} and \texttt{Organization}, since its \ac{kr} (derived from documents metadata) of reference is limited to these two types.}
    % (``Implicit'' derives from RAE, since rule-based additional extractors perform recognition and linking as a whole)
    \label{tab:algotypes}
\end{table}

%inally, we observe that \ac{nel} and NIL clustering can be viewed as two faces of the same process aimed at consolidating the mentions found by \ac{ner} algorithms into a \textit{knowledge layer}: mentions linked to the same knowledge base entity are, in fact, clustered together, while NIL mentions clustered together are linked to the same new entity. As discussed in \Cref{sec:related_work}, this process of knowledge consolidation is the most novel aspect in our proposal, the one whose feasibility we want to evaluate and the one we mainly focus on in this paper.

When multiple \ac{ner} algorithms are employed, it is necessary to combine their extracted annotations and resolve any resulting conflicts. To this end, we introduce heuristic-based combination rules. Both the \ac{ner} algorithms and the combination rules adopted in this work are described in detail in the remainder of this section.






% TODO SCHEMA

% TODO pipeline

\subsubsection{Entity Recognition}          
For the \ac{ner} task, we selected two general-domain, ML-based algorithms for Italian: \spacy~\cite{honnibal2020spacy} and Tint~\cite{2018tint2}. 
In addition, we developed a rule-based service, referred to as ``Trie\ac{ner}'', which performs efficient gazetteer-based lookup using a prefix tree.  
Each algorithm recognizes a different set of entity types. \Cref{tab:algotypes} reports the correspondence between the entity types in our benchmark dataset and those detected by each algorithm, including their combination.  

\ac{ner} with \spacy is carried out using the Italian transformer model ``dbmdz/bert-base-italian-uncased''\footnote{\url{https://huggingface.co/dbmdz/bert-base-italian-uncased}}, pretrained by \textcite{stefan_schweter_2020_4263142}.  
The model produces contextual representations that are then processed by a neural transition-based parser (see \Cref{sec:rw:ner}).  
We fine-tuned the \spacy pipeline---including both the transformer and the parser---on the Italian WikiNER dataset~\cite{NOTHMAN2013151}, which is based on Wikipedia articles annotated with \texttt{Person}, \texttt{Location}, \texttt{Organization}, and \texttt{Miscellaneous} entities.

Tint~\cite{2018tint2} is an Italian NLP suite based on Stanford CoreNLP~\cite{manning-EtAl:2014:P14-5}.  
Its \ac{ner} module combines CRF-based sequence taggers~\cite{Lafferty2001} with rule-based recognizers for \texttt{Money}, \texttt{Number}, and temporal expressions (\texttt{Time}).

Trie\ac{ner} relies on a gazetteer efficiently indexed with a prefix tree (trie)~\cite{s-yata-marisa-trie}, following prior work on tree-based \ac{ner}~\cite{Nguyen_tree_ner_2014}.  
Our gazetteer is built from the document metadata: entity labels are tokenized, enriched with all name permutations (for personal names), and inserted into the trie.  
The algorithm supports partial matches---\eg detecting ``Smith'' when the full label ``Jane Smith'' exists in the gazetteer.  
Unlike standard \ac{ner}, Trie\ac{ner} also provides the link between the matched text and the originating entity (or entities, if shared).  
This approach leverages metadata fields such as plaintiffs, defendants, and judges to increase the likelihood that these entities are also recognized in the text.

Suppose ``Jane Smith'' is listed as a plaintiff in the metadata: Trie\ac{ner} searches all token permutations and detects both ``Smith'' (partial match) and ``Smith Jane''.

Finally, we merge the annotations produced by the different algorithms.  
Each \ac{ner} annotation consists of a text span associated with a type, so conflicts may arise when spans overlap or when overlapping spans have inconsistent types---for example:
\begin{quote}
``\underline{Mr. Cityville}[PER]'' (from one algorithm) vs. ``\underline{Cityville}[LOC]'' (from another),  
or ``Italy'' labeled both as \texttt{Location} and \texttt{Organization}.
\end{quote}

We resolve these conflicts as follows:
\begin{enumerate}
    \item Span conflicts: overlapping annotations can be total (same boundaries) or partial.  
    In partial overlaps, we select the longest annotation, assuming it provides the most information.  
    As an exception, for entities of type \texttt{Person}, we deprioritize spans longer than $k$ tokens, using $k=6$ after identifying it as the maximum number of tokens for names in ICCJ900k metadata.
    \item Type conflicts: different algorithms may predict different types.  
    To handle this, we assign a weight to each algorithm and perform a \textit{weighted majority vote} to select the final type.  
    If multiple algorithms predict the same type, their weights are summed.  
    This mechanism allows us to control each algorithm's influence on the final decision.
\end{enumerate}



\subsubsection{Knowledge consolidation}
\label{sec:claw_consolidation}
% We opted for the BLINK bi-encoder because it provides a semantic representation of mentions and entities enabling us to represent newly identified entities starting from their mentions so that it is possible to use them as target entities for \ac{nel}. % TODO migliorare
% Provato a scriverla meglio
% As in \acl{inel}, we chose the bi-encoder paradigm because it is able to produce a semantic representation of both mentions and entities. This property allows us to calculate the semantic similarity of entity mentions during the NIL clustering step and to represent new entities using their mentions so that the \ac{nel} system is able to link to the new entities~\cite{ijckg2022}.
The knowledge consolidation process comprises three tasks: \ac{nel}, NIL prediction, and NIL clustering.  
The pipeline is adapted from the English version introduced in \textcite{ijckg2022} and detailed in \Cref{sec:ijckg:pipeline}. It consists of the BLINK bi-encoder~\cite{blink}, a logistic regression classifier for NIL prediction, and a three-step algorithm for NIL clustering.  
Although we do not evaluate an incremental scenario with legal documents, we adopt a bi-encoder-based architecture because it supports the representation of new entities---a crucial property for the legal domain.

To enable \ac{nel} for the Italian language, we trained a bi-encoder following the methodology proposed by \textcite{blink} for English.  
We initialized the encoder with weights from the Italian BERT-base model ``dbmdz/bert-base-italian-uncased''\footnote{\url{https://huggingface.co/dbmdz/bert-base-italian-uncased}}~\cite{stefan_schweter_2020_4263142} and fine-tuned it on 9 million mention--entity pairs derived from Italian Wikipedia hyperlinks.  
Training was conducted for four epochs with in-batch random negatives, using the AdamW optimizer~\cite{loshchilov2018decoupled} with an initial learning rate of $1 \times 10^{-5}$ and a batch size of 20.  
A fifth epoch was then performed using hard negatives---one per sample---selected as the incorrect entity with the highest linking score for the given mention.  
As the \acl{kr} we use the $\approx$1.5M entities extracted from Italian Wikipedia\footnote{\url{https://it.wikipedia.org}} after filtering out redirects and disambiguation pages.  
In the remainder of this work, we refer to the resulting model as \blinkita.

For NIL prediction, we use the same Wikipedia-based dataset employed to train the \ac{nel} model.  
As in \Cref{sec:ijckg:pipeline}, we adopt a \textit{logistic regression classifier} that takes as input:
\begin{enumerate}
    \item the \ac{nel} score of the top-ranked entity; and
    \item the difference between the top score and that of the second-best candidate (\textit{secondiff}).
\end{enumerate}
The classifier outputs a probability $p \in [0,1]$, where $p = 1$ indicates that the top-ranked entity is likely correct for the given mention, and $p = 0$ indicates the opposite.  
In the latter case, the mention is treated as NIL, under the assumption that if the correct entity is not top-ranked, it is not present in the \ac{kr}.

Finally, NIL clustering follows the three-step procedure described in \Cref{sec:ijckg:pipeline}.  
The clustering thresholds are determined through grid search on the same Wikipedia-based dataset used for \ac{nel} training.


% spostare negli experiments
% We conduct a first evaluation of \blinkita on two italian datasets: the tweet-based NEEL-IT@Evalita 2016~\cite{basile2016neelit} and the news-based VoxEL~\cite{rosales2018voxel}. The results are available in \Cref{tab:nereval:atomic}.

% TODO voxel multilingual \ac{nel} dataset

% Although a systematic comparison is not possible (as further discussed in ~\Cref{app:blink}), our goal is to have some insights on the performance of \blinkita in a challenging linking scenario~\cite{basile2016neelit}.

% The insights suggest that \blinkita achieves reasonable performance, despite some known limitations of its current implementation (see~\Cref{app:blink}).
% todo ?

%the 
%Unfortunately a systematic comparison is not possible due to several reasons has been evaluated on the dataset NEEL-IT@Evalita 2016~\cite{basile2016neelit}, that is composed of tweets in the Italian language; see \Cref{tab:blink_ita_eval}.


% The NIL clustering sub-service is a three-step algorithm~\cite{simone_24}: initially, the mentions are clustered using their surface form (\ie edit-distance $\leq 3$ for words longer than 3 characters, perfect match otherwise), then inside each cluster we apply a hierarchical clustering algorithm by using a predefined threshold to split clusters, considering the mentions' dense vectors provided by the bi-encoder that correspond to a semantic representation of the mention.
% In this way, each cluster from the first step is divided into sub-clusters.
% Finally, semantically similar sub-clusters are merged together, according to the distance between their medoid vectors.

% The NIL prediction and NIL clustering approaches are based on incremental extraction methods from English documents~\cite{ijckg2022}.


% Per stabilire quale span di testo preferire, nel caso in cui si tratti di persona
% o categorie correlate ad essa, si sceglie l'annotazione più lunga, a patto che
% essa sia composta da un numero di token inferiore o uguale al massimo
% determinato per via empirica nel corpus (6 elementi nel caso in esame), nel
% caso di altre categorie o qualora non sia possibile stabilire una categoria,
% viene preferita l'annotazione più lunga.




\subsection{Evaluating a General-domain Pipeline on Italian Civil Judgments}
\label{sec:judgment_ee_eval}

We evaluate the entity extraction pipeline in the following experimental settings:
\begin{enumerate}
    \item \textit{Atomic}: we analyze the performance of each algorithm atomically to isolate it from the effects of error propagation. In practice, before analyzing the performance of each task, the prediction from the previous tasks in the pipeline are corrected according to the ground truth. For \ac{ner} we use the 146-documents ICCJ146-EE as the evaluation dataset and the 30-documents subset ICCJ30-\ac{inel} for the knowledge consolidation tasks. \ac{nel}, NIL prediction, and NIL clustering are applied, and evaluated, only on annotations of type \texttt{Person}, \texttt{Location}, \texttt{Organization}. This analysis aligns with the concept of \acf{hitl} validation, in which the results of each step are refined by humans prior to being utilized as input for the subsequent stages within the pipeline. For \ac{ner}, we additionally calculate performance metrics on a per-type basis and with different matching criteria, including more relaxed ones that also accept spans that are partially correct (defined below).  
    \item \textit{\ac{nel} with NIL prediction}: we study the combined performance of \ac{nel} and NIL prediction, since these two tasks are closely related, isolating them from \ac{ner} errors. To achieve this, we execute \ac{nel} and NIL prediction on ground truth \ac{ner} annotations from ICCJ30-\ac{inel}.
    \item \textit{End-to-end}: we evaluate the performance of our system on each mention in ICCJ30-\ac{inel} considering \ac{ner}, \ac{nel}, and NIL prediction---as done, for example, in the NEEL challenge~\cite{basile2016neelit}. In this setting, each mention can be either linked to an entity in the \ac{kr} or classified as NIL. We exclude NIL clustering from this evaluation, since assessing whether a NIL mention is in the correct cluster would overly complicate the analysis; instead, clustering is evaluated separately in the \textit{atomic} setting.


\end{enumerate}
To better study the difficulty of the domain, we evaluate our \ac{ner} and \ac{nel} systems also on standard benchmarks in the Italian language: WikiNER~\cite{NOTHMAN2013151} and I-CAB~\cite{magnini-etal-2006-cab} for \ac{ner}; Italian VoxEL~\cite{rosales2018voxel} and NEEL-IT~\cite{basile2016neelit} for \ac{nel}.

In the remainder of this section, we present the evaluation metrics, the matching criteria for \ac{ner}, and the results.

% The following sections contain the explanation of the evaluation metrics (\Cref{sec:metrics})
% %, the description of the ground truth (\Cref{}),
% and the results of the experiments (\Cref{sec:results}).

% %%%


% The experiments were carried out to evaluate the algorithms with ground truth labels, focusing on \ac{ner}, \ac{nel}, NIL prediction, and NIL clustering. 
% The aim of this evaluation is to analyze the difficulty of \ac{ner}, \ac{nel}, NIL prediction, and NIL clustering algorithms without any domain-specific fine-tuning. We also seek to establish a best-effort baseline, which can be further improved by adapting the algorithms to the specific target domains, as demonstrated in the existing literature under similar contexts with limited manually annotated data available~\cite{ehrmann2021named}.

%%%%


% The experiments carried out for the evaluation of the \textit{main pipeline} aim to investigate the feasibility of applying the full entity extraction pipeline in a specific domain as the one characterized by Italian court judgments. By evaluating the performance of \ac{ner}, \ac{nel}, NIL prediction, and NIL clustering algorithms without additional fine-tuning, we want to analyze the difficulty of each task and of the whole entity extraction task. Also, we intend to establish a best-effort baseline, which can be improved by adapting algorithms to the target domains, by using approaches proposed in similar contexts where only a limited number of manually annotated data can be collected~\cite{ehrmann2021named}.  % ner in contesti storici
%
% By independently analyzing each step of the pipeline we also study the propagation of errors between steps.
% %, which can lead to inaccurate results and make it difficult to identify the cause of errors. 
% The analysis of individual steps also fits the proposed \ac{hitl} approach since the input of each step can be obtained through a human refinement of the algorithm's output of the previous step.
%
% By conducting independent analyses of each step in the pipeline, we can study the propagation of errors between these steps.
% These analyses align with the concept of \ac{hitl} validation, in which the results of each step can be refined by humans prior to being utilized as input for the subsequent stages within the pipeline.
%
% The types involved in each step of the experimentation are listed in \Cref{tab:algotypes}.

\subsubsection{Evaluation Metrics and Criteria}\label{sec:metrics}
% \paragraph{\ac{ner}}
For the \textit{atomic} evaluation of \ac{ner}, we calculate precision, recall, and F$_1$ score with different matching criteria that determine when an extracted mention is correct with respect to the human-annotated ground truth.
Indeed, in some contexts, ``partially-correct'' annotations might be acceptable---for instance, identifying ``A4 Highway'' instead of the more specific ``A4 Highway Torino-Trieste'' might be informative enough in some settings~\cite{nadeau2007survey, tsai2006various}. Furthermore, in a \ac{hitl} scenario, we prefer to detect partially-correct annotations rather than to miss them, since a human can quickly correct annotation boundaries.

% The metrics used for the evaluation of \ac{ner} are inspired by~\cite{tsai2006various}, and have different criteria to mark a prediction as correct. In this way, it is possible to compute precision, recall, and F$_1$ score at different severity levels, thus providing a better overview of the effectiveness of the algorithm.

The adopted criteria---taken from~\cite{tsai2006various} and extended to provide both a typed and untyped evaluation---allow us to assess the behavior of the algorithms at different severity levels. The metrics are defined as follows:
\begin{enumerate}
    \item \textit{Strong:} the predicted entity has an exact span match with the ground truth annotations.
    \item \textit{Strong-typed:} the predicted entity has an exact span match and the type is correct with respect to the ground truth annotations.
    \item \textit{Approximate:} the predicted entity is contained in the correct span (or vice versa).
    \item \textit{Approximate-typed:} the predicted entity is contained in the correct span (or vice versa) and the type is correct.
    \item \textit{Partial:} the predicted entity is overlapping with the correct span.
    \item \textit{Partial-typed:} the predicted entity is overlapping with the correct span and the type is correct.
\end{enumerate}

% \todo{controlla HITL}

For the sake of clarity, let $c$ denote the chosen evaluation criterion, $Y^i_c$ the number of correctly predicted annotations according to $c$, $Y^i$ the total number of predicted annotations, and $Y_{GT}$ the number of annotations in the ground truth.  
Precision, recall, and F$_1$ score are then defined as:
\begin{equation}
P = \frac{Y^i_c}{Y^i}, \qquad
R = \frac{Y^i_c}{Y_{GT}}, \qquad
F_1 = 2 \times \frac{P \times R}{P + R}.
\end{equation}
Here, $P$ measures the fraction of correct predictions, $R$ measures coverage over the ground truth, and F$_1$ represents their harmonic mean.



% \paragraph{Knowledge consolidation} % (\ac{nel}, NIL prediction, NIL clustering)
% The evaluation of \ac{nel} is composed of two parts. The first part regards the base entity linking, which is evaluated on \textbf{recall@K} metric to assess if the correct link is contained in the top-K candidates with the highest score. The second part evaluated is the NIL predictor. Since it is a binary classifier, it is evaluated on the standard metrics of precision, recall, and F$_1$ score. 
% The knowledge consolidation tasks are evaluated independently from \ac{ner} and also independently from each other.

% Each consolidation step, namely \ac{nel}, NIL prediction, and NIL clustering, has been evaluated atomically to exclude the effects of error propagation. This evaluation setting is also in line with the scenario in which a \ac{hitl} validates the output of each step.

% \todo{valutare se la figura si può togliere}
% \Cref{fig:expected_outcome} shows the expected system behavior given a mention $m$ when it refers to an entity that exists in the \ac{kr} and when it is NIL.

% \paragraph{\ac{nel}}
\ac{nel} is \textit{atomically} evaluated on accuracy, that is the number of mentions linked to the correct entity divided by the total number of mentions to link, and recall@k. This latter metric assesses whether the correct entity is among the top-k candidates with the highest linking score (in our setting, accuracy and recall@1 are equivalent). Note that \ac{nel}, and also NIL prediction and NIL clustering, is evaluated solely on the types \texttt{Person}, \texttt{Location}, \texttt{Organization}. 

\ac{nel} is evaluated \textit{atomically} using accuracy and recall@k.  

\begin{equation}
\text{Accuracy} = \frac{\text{Number of correctly linked mentions}}{\text{Total number of mentions to link}},
\end{equation}

\begin{equation}
\text{Recall@k} = \frac{\text{Number of mentions for which the correct entity is among the top-}k\text{ candidates}}{\text{Total number of mentions to link}}.
\end{equation}




% \paragraph{\ac{nel}}
% \paragraph{NIL prediction}
The NIL prediction classifier is first evaluated \textit{atomically}, computing precision, recall, and F$_1$ for both the NIL and $\neg$NIL classes. Since NIL prediction is tightly coupled with \ac{nel}, we decouple it by counting as correct also the cases where a mention linked to a wrong entity is classified as NIL: in such cases the top-ranked entity according to \ac{nel} is incorrect, and the NIL classifier correctly assumes that the true entity is not present in the \ac{kr}.

NIL clustering is also evaluated \textit{atomically} using MUC, B$_3$, and CEAF$_e$ as in \Cref{sec:inel_metrics}, reporting precision, recall, and F$_1$ for each. To avoid error propagation, in this setting, NIL clustering is run only on the NIL mentions from the ground truth.

Finally, we conduct a joint evaluation of \ac{nel} with NIL prediction. Starting from the reference \ac{ner} annotations and following \Cref{sec:ijckg:evaluation_procedure}, we compute the accuracy on:
\begin{enumerate}
\item mentions that should be linked to the \ac{kr}---as in (a) of \Cref{sec:ijckg:evaluation_procedure},
\item mentions that should be classified as NIL---as in (b),
\item all mentions---as in (d).
\end{enumerate}
These correspond, respectively, to the fraction of correctly processed $\neg$NIL mentions, NIL mentions, and all mentions.

% \begin{figure}
%     \centering
%     \includegraphics[
%     width=.4\textwidth,
%     keepaspectratio]
%     {experiments_expected.drawio.png}
%     \caption{Schema of the expected outcome of the consolidation tasks, given a mention m
% referring to an entity E, when (a) E is known ``a priori'' in the \ac{kr} or (b) E is not in the
% \ac{kr} (NIL).}
%     \label{fig:expected_outcome}
% \end{figure}

% Each task has been atomically evaluated to avoid error propagation and with the assumption that a \ac{hitl} validated the output from the previous step.
% Additionally, we provide an evaluation of the entire consolidation pipeline by calculating the accuracy of the system behavior starting from a human-validated \ac{ner} for all the input mentions. We also provide distinct results for the mentions that should be linked both to \ac{kr} and out-of-\ac{kr} (\ie NIL) entities. \Cref{fig:expected_outcome} shows the expected system behavior given a mention $m$ when it refers to an entity that exists in the \ac{kr} and when it is NIL.
% Additionally, for the evaluation of the entire consolidation pipeline, independently by the \ac{ner}, we calculate the accuracy of the system behavior for all the input mentions, and separately for the mentions that should be linked to the \ac{kr} and for the ones that should be classified as NIL. \Cref{fig:expected_outcome} shows the expected system behavior of the system given a mention $m$ when it refers to an entity that exists in the \ac{kr} and when it is NIL.

Finally, for the \textit{end-to-end} evaluation of \ac{ner}, \ac{nel}, and NIL prediction, we define two extended criteria: \textit{approximate linking} and \textit{approximate-typed linking}. These criteria consider an annotation correct only if its \ac{ner} boundaries---and \ac{ner} type, for typed linking---are correct, it is linked to the correct entity and classified as $\neg$NIL or correctly classified as NIL. For both criteria, we report precision, recall, and F$_1$ score.


% \subsubsection{Scalability, Extensibiltiy, and Availability}
% Finally, we study the scalability of the proposed architecture by calculating the required time for the execution of the generic algorithms and of the additional rule-based extractors, which recognize postal addresses and law articles.

%\todo{la teniamo questa considerazione sui processi asincroni commentata? nella discussione?}
% We argue that several operations are performed by independent systems that can run asynchronously in any order (apart for workflows), as they produce independent annotations. These operations include the \ac{ner} algorithms and the rule-based additional extractors, while other services like \ac{nel} strictly depend on the previous steps.

The hyperparameters for the combination rules were set prioritizing Trie\ac{ner}, followed by \spacy, and then Tint. This choice reflects the reliability of Trie\ac{ner}, which leverages in-domain knowledge from metadata, and the fact that \spacy employs more recent models than Tint. The resulting weights are 0.6 for Trie\ac{ner}, 0.3 for \spacy, and 0.1 for Tint, ensuring that in cases of type conflict, Trie\ac{ner}'s prediction---based on metadata---takes precedence.


% \paragraph{Overall}


\subsection{Results}\label{sec:results}

% \subsubsection{\ac{ner}}
% Since the actual experimentation does not involve algorithms for class specialization, which is left for future works, the types of \textit{Judge}, \textit{Lawyer}, \textit{Court}, \textit{Plaintiff}, \textit{Defendant}, and \textit{Law Article} have been excluded from the evaluation.

\begin{table}[H]
    \centering
    \begin{tabular}{ll|rrr}
        \toprule
        \textbf{Model} & \textbf{Dataset} & \textbf{Precision} & \textbf{Recall} & \textbf{F$_1$} \\
        \midrule
        \spacy & WikiNER~\cite{NOTHMAN2013151} & 91.9 & 91.9 & 91.9 \\
        GilBERTo~\cite{gilberto} & WikiNER~\cite{NOTHMAN2013151} & 92.7 & 92.7 &	92.8 \\
        BERTino~\cite{muffo2023bertino} & WikiNER~\cite{NOTHMAN2013151} & - & - & 90.4 \\
        BERTino Teacher Model~\cite{muffo2023bertino} & WikiNER~\cite{NOTHMAN2013151} & - & - & 91.8 \\
        Tint & I-CAB~\cite{magnini-etal-2006-cab} & 84.4 & 80.0 & 82.1 \\
        \bottomrule
    \end{tabular}
    \caption{Results obtained by our algorithms (\spacy and Tint) on Italian benchmark datasets compared with recent competitive approaches. I-CAB~\cite{magnini-etal-2006-cab} contains the types \texttt{Person}, \texttt{Location}, and \texttt{Organization}. WikiNER~\cite{NOTHMAN2013151} additionally considers the type \texttt{Miscellaneous}.}
    \label{tab:sotaner}
\end{table}

First, \Cref{tab:sotaner} reports a comparison between two of our \ac{ner} algorithms, \spacy and Tint, on public benchmark datasets. For reference, we also include the results of two recent competitive approaches: GilBERTo~\cite{gilberto} and BERTino~\cite{muffo2023bertino}. As shown, our \spacy model, which leverages a BERT encoder, performs on par with state-of-the-art systems.

\Cref{tab:nereval:bymetric} presents the results of the \textit{atomic} evaluation of \ac{ner}. As expected, there is a clear gap between the scores obtained under the \textit{strong} and \textit{approximate} criteria, indicating that while our system often detects entities correctly, it sometimes fails to identify their exact boundaries. Such boundary errors can, however, be corrected efficiently through \ac{hitl} revision. The difference between the \textit{approximate} and \textit{partial} criteria is negligible; therefore, we regard the \textit{approximate} criterion as sufficiently permissive and omit the \textit{partial} one in the remainder of this evaluation.

\begin{table}[H]
    \centering
    \begin{tabular}{l|rrr|rrr}
        \toprule
        & \multicolumn{3}{c|}{\textbf{Untyped}} & \multicolumn{3}{c}{\textbf{Typed}} \\
        & \textbf{Precision} & \textbf{Recall} & \textbf{F$_1$} & \textbf{Precision} & \textbf{Recall} & \textbf{F$_1$} \\
        \midrule
        % approximate typed % data denaro luogo organizzazione persona, ( con miscellaneous )
        Strong              & 51.9 &	43.7 &	47.5 & 44.7 &	37.7 &	40.9  \\
        Approximate         & 86.9 &	68.0 &	76.3 & 64.6 &	53.1 &	58.3 \\
        Partial             & 87.4 &	68.6 &	76.9 & 64.9 &	53.3 &	58.6 \\ 
        \bottomrule
    \end{tabular}
    \caption{\textit{Atomic} \ac{ner} evaluation (precision, recall, F$_1$ score) with different matching criteria (rows), typed and untyped. These results are obtained using the combination of the \ac{ner} algorithms and considering all the evaluated types.}
    \label{tab:nereval:bymetric}
\end{table}

In \Cref{tab:nereval:bytype}, we provide a per-type comparison of the \textit{atomic} evaluation for both the individual \ac{ner} algorithms and their combination. The entity types \texttt{Date}, \texttt{Money}, and \texttt{Miscellaneous} are each handled by a single algorithm---Tint for the first two, and \spacy for the latter. The seemingly counterintuitive differences between the combined results and those of the algorithms responsible for these types arise from boundary conflicts, which may also occur between mentions of different types.

\Spacy proves particularly effective in identifying \texttt{Person} and \texttt{Location} entities, whereas Tint performs better on \texttt{Organization}. The Trie\ac{ner} evaluation yields high precision but low recall, as it can only recognize entities present in its dictionary, which mainly includes persons and a few organizations.

\begin{table}[H]
    \centering
    \begin{tabular}{l|l|rrr}
        \toprule
         \textbf{Type}         & \textbf{Algorithm} &   \textbf{Precision}   &   \textbf{Recall}   &   \textbf{F$_1$} \\
        \midrule
        \texttt{Person}        & \spacy   & \textbf{92.1} & 76.3 & \textbf{83.5}\\
              & Tint    & 90.3 & 62.1 & 73.6\\
              & Trie\ac{ner}    & 76.7 & 34.6 & 47.7\\
              & Combination  & 81.5 & \textbf{79.9} & 80.7\\
        \midrule
        \texttt{Location}      & \spacy   & 35.4 & \textbf{90.8} & 51.0\\
                      & Tint    & 49.2 & 82.6 & \textbf{61.7}\\
                      & Trie\ac{ner}    & - & - & -\\
                      & Combination  & \textbf{59.9} & 60.1 & 60.0\\
        \midrule
        \texttt{Organization}  & \spacy   & 58.2 & 40.9 & 48.1\\
                      & Tint    & 65.6 & 56.5 & \textbf{60.7}\\
                      & Trie\ac{ner}    & \textbf{90.0} & 0.2 & 0.5\\ % TODO togliere: we omit trie org because the number of org is not statistically relevant
                      & Combination  & 34.2 & \textbf{92.0} & 49.9\\
        \midrule
        \texttt{Date}          & \spacy   & - & - & -\\
                      & Tint    & \textbf{83.7} & \textbf{55.2} & 66.5\\
                      & Trie\ac{ner}    & - & - & -\\
                      & Combination  & \textbf{83.7} & 55.1 & 66.5\\
        \midrule
        \texttt{Money}         & \spacy   & - & - & -\\
                      & Tint    & \textbf{98.1} & \textbf{57.3} & \textbf{72.3}\\
                      & Trie\ac{ner}    & - & - & -\\
                      & Combination  & \textbf{98.1} & 56.9 & 72.0\\
        \midrule
        \texttt{Miscellaneous} & \spacy   & 25.1 & \textbf{4.1} & \textbf{7.0}\\
                      & Tint    & - & - & -\\
                      & Trie\ac{ner}    & - & - & -\\
                      & Combination  & \textbf{26.4} & 3.9 & 6.8 \\
        \midrule  \textbf{Overall*} & \spacy   & 63.8 & 43.2 & 51.5\\
                      & Tint    & 74.8 & 60.1 & 66.6 \\
                      & Trie\ac{ner}    & \textbf{76.7} & 11.3 & 19.7 \\
                      & Combination  & 66.0 & \textbf{67.6} & \textbf{66.8} \\
        \bottomrule
        % \midrule
        % \multicolumn{5}{c}{Additional Rule-based Extractors} \\
        % \midrule
        % Postal Address       & rule-based & .902 & .881 & .892 \\
        % Law Articles  & rule-based & 1.0 & .702 & .825 \\
        % \midrule
        % \multicolumn{5}{c}{Benchmark Comparison**} \\
        % \midrule
        % WikiNER~\cite{NOTHMAN2013151} & \spacy & .919 & .919 & .919 \\ % per org loc misc
        % I-CAB~\cite{magnini-etal-2006-cab} & Tint & .844 & .800 & .821 \\ % per org loc
    \end{tabular}
    \caption{\textit{Atomic} \ac{ner} evaluation by type and algorithm using \textit{approximate-typed match}. *\texttt{Miscellanea} mentions are excluded from the overall calculation. **Note that the type Organization is underrepresented in the \acl{kr} used by Trie\ac{ner}, thus only a few entities of this type (\ie $\approx$10) have been predicted.}
    % TODO Trie\ac{ner}
    \label{tab:nereval:bytype}
\end{table}

Our combination rules offer a good balance between precision and recall, achieving an F$_1$ score of 66.8% under the \textit{approximate matching} criterion, even though they do not surpass the best-performing individual algorithm for each type.

The results of \spacy and Tint on standard benchmarks confirm the challenging and heterogeneous nature of our target domain. Comparing \Cref{tab:sotaner,tab:nereval:bymetric}, we observe a substantial performance drop when applying \ac{ner} models to domain-specific data. A closer look at the per-type evaluation (\Cref{tab:nereval:bytype}) shows that performance generally decreases across most categories compared to benchmarks. Notably, the \texttt{Person} class maintains relatively strong results (F$_1 > 80\%$), while the recognition of \texttt{Miscellanea} entities is unsatisfactory---likely due to the distinct nature of domain-specific \texttt{Miscellanea} compared to those in standard corpora. For the remaining categories (\texttt{Location}, \texttt{Organization}, \texttt{Date}, and \texttt{Money}), F$_1$ scores range between 50\% and 70\%, suggesting that although the general-domain models can handle these types to some extent, domain-specific fine-tuning could yield substantial gains.

% \begin{table}[H]
%     \centering
%     \begin{tabular}{l|lll|lll}
%         & \multicolumn{3}{l|}{strong-typed} & \multicolumn{3}{l}{approximate typed} \\
%         \midrule
%         & P & R & F$_1$ & P & R & F$_1$ \\
%         \hline
%         % approximate typed % data denaro luogo organizzazione persona, no altro
%         \spacy  & .490 & .339 & .401 & .638 & .432 & .515 \\
%         Tint   & \textbf{.511} & .417 & .459 & .748 & .601 & .666 \\
%         Trie\ac{ner}   & .174 & .026 & .045 & \textbf{.767} & .113 & .197 \\
%         Combination & .461 & \textbf{.488} & \textbf{.474} & .660 & \textbf{.676} & \textbf{.668} \\
%         % \midrule
%         % Macro avg & P & R & F$_1$ & P & R & F$_1$ \\
%         % \hline
%         % \spacy &&&&&& \\
%         % Tint &&&&&& \\
%         % Trie\ac{ner} &&&&&& \\
%         % Combination &&&&&& \\
%        % Overall  & .919 & .919 & .919 \\
%        % Person  & .969 & .967 & .968 \\
%        % Location  & .913 & .922 & .918 \\
%        % Organization  & .873 & .920 & .896 \\
%        % Miscellaneous  & .864 & .823 & .843
%     \end{tabular}
%     \caption{Comparison of the \ac{ner} algorithms, including their combination, according to the strong-typed and approximate typed criteria. The values correspond to the micro-averaged precision, recall, and F$_1$ score of the types ``person'', ``location'', ``organization'', ``money'', and ``date'' (``miscellaneous'' is excluded).}
%     \label{tab:nereval}
% \end{table}


% \subsubsection{\ac{nel} (entity linking + NIL Prediction)}
% \subsubsection{Knowledge consolidation (Entity Linking + NIL prediction + NIL clustering)}
% (although which corresponds to the percentage of mentions whose correct entity for linking is among the top 100, suggesting that a re-ranking phase (\eg a cross-encoder \cite{blink}) could be beneficial.

The results of the atomic evaluation for the various knowledge consolidation tasks, as well as the joint \ac{nel} and NIL prediction evaluation, are reported in \Cref{tab:nereval:atomic}. The \ac{nel} module performs satisfactorily, achieving an accuracy of 73.5\% and a recall@100 of 90.8\%, consistent with results obtained on benchmark datasets which are available in \Cref{tab:nelsota}.

Regarding NIL prediction, the classifier effectively identifies NIL mentions but struggles with $\neg$NIL ones, yielding an overall F$_1$ score of 64.5\%. The main weakness of the consolidation pipeline lies in the NIL clustering stage, which achieves low F$_1$ scores under both B$^3$ and CEAF$_e$ metrics.

When jointly evaluating \ac{nel} and NIL prediction (disregarding \ac{ner} errors), 79.1\% of mentions are correctly processed. Among these, NIL mentions are identified with very high accuracy (91.9\%), whereas mentions that should be linked to the \ac{kr} reach only 46.8\% accuracy. This discrepancy can be attributed to a bias in the NIL classifier, which tends to over-predict the NIL class.

\begin{table}[H]
    \centering
    \begin{tabular}{lr|rrrr}
        % \multicolumn{2}{c|}{\textbf{Task}} & \multicolumn{3}{c}{\textbf{Metric}} \\
        \toprule
        \textbf{\Acl{nel}} & Accuracy & 73.5 & & & \\
        & Recall@100 & 90.8 & & & \\
        \midrule
        \textbf{NIL prediction} & & \textbf{Precision} & \textbf{Recall} & \textbf{F$_1$} \\
        & NIL & 92.2 & 86.5 & 89.2 \\
        & $\neg$NIL & 58.5 & 72.0 & 64.5 \\
        \midrule
        \textbf{NIL clustering} & & \textbf{Precision} & \textbf{Recall} & \textbf{F$_1$} \\
        & MUC & 71.9 & 83.9 & 77.4 \\
        & B$^{3}$ & 16.4 & 60.1 & 25.8 \\
        & CEAF$_{e}$ & 7.2 & 31.7 & 11.7 \\
        \midrule
        % \multicolumn{1}{l|}{\textbf{\ac{nel} \& NIL prediction}} & Accuracy & & \\
        % in \ac{kr}     & .468  &   &   \\ 
        % NIL    & .919  &   &   \\
        % All               & .791  &   &   \\
        \textbf{\Ac{nel} \& NIL prediction} & & \textbf{in \ac{kr}} & \textbf{NIL} & \textbf{All} \\
        & Accuracy & 46.8 & 91.9 & 79.1 \\
        \bottomrule
    \end{tabular}
    \caption{\textit{Atomic} evaluation of the knowledge consolidation tasks and joint evaluation of \ac{nel} and NIL prediction.}
    \label{tab:nereval:atomic}
\end{table}

\begin{table}[h]
    \centering
    \begin{tabular}{l|rrr}
        % \multicolumn{2}{c|}{\textbf{Task}} & \multicolumn{3}{c}{\textbf{Metric}} \\
        \toprule
        & \textbf{s-VoxEL-it} & \textbf{r-VoxEL-it} & \textbf{NEEL-IT} \\
        \midrule
        \textbf{Accuracy}  & 88.9 & 64.7 & 69.0 \\
        \textbf{Recall@100} & 96.8 & 91.5 & - \\
        \bottomrule
    \end{tabular}
    \caption{\ac{nel} results on benchmark datasets: strict (s-) and relaxed (r-) version of Italian VoxEL~\cite{rosales2018voxel} and NEEL-IT@Evalita 2016~\cite{basile2016neelit}. for NEEL-IT, Twitter profile tags (@username) and hashtags (\#tag) were filtered out.}
    \label{tab:nelsota}
\end{table}

% overall
Finally, \Cref{tab:nereval:withlink} presents the end-to-end evaluation results using the \textit{approximate linking} and \textit{approximate-typed linking} criteria. The \textit{approximate linking} setup, which ignores the type predicted with \ac{ner}, yields noticeably better performance than \textit{approximate-typed linking}: the F$_1$ score reaches 64.0\%, nearly eight points higher. This result suggests that once a mention is successfully linked by \ac{nel}, the linked entity can help refine or even correct the type initially assigned by \ac{ner}.
% todo we leave this to future works


\begin{table}[H]
    \centering
    \begin{tabular}{lrrr}
        \toprule
        & \textbf{Precision} & \textbf{Recall} & \textbf{F$_1$} \\
        \midrule
        % approximate-typed % luogo organizzazione persona, no altro data denaro 
        Approximate linking      & 59.8 & 68.9 & 64.0 \\
        Approximate-typed linking & 52.3 & 60.9 & 56.3 \\
        \bottomrule
    \end{tabular}
    \caption{\textit{end-to-end} \ac{ner}, \ac{nel}, and NIL prediction evaluation.}
    \label{tab:nereval:withlink}
\end{table}

\subsection{Discussion}
In the \textit{end-to-end} evaluation, the system achieved an F$_1$ score of 64.0\% (under the \textit{approximate linking} criterion), meaning that more than one mention out of two is correctly processed across \ac{ner}, \ac{nel}, and NIL prediction. Considering the domain's high specificity, the noisy input, and the absence of in-domain training or fine-tuning, these results represent a promising baseline for entity identification and linking.

The main weakness lies in the final NIL clustering stage, which obtained low B$^3$ and CEAF$_e$ scores. Interpreting clustering metrics is inherently difficult, and prior work has shown wide variability across datasets and evaluation settings~\cite{streaming_crossdoc_coreference}. Since suboptimal clustering directly affects the consolidation of NIL mentions, improving this step is crucial.
% A possible direction is to tune clustering hyperparameters (\eg thresholds) on domain-specific data, as suggested by recent studies~\cite{streaming_crossdoc_coreference}.\todo{which params we used?}

Also, the NIL prediction component shows a bias toward the NIL class, which could be reduced through domain-specific training. However, in the context of court judgments, this limitation is less problematic, as most key entities (\eg plaintiffs, defendants, judges, and attorneys) are in fact often NIL.

Performance degradation was also observed in \ac{ner} models trained on publicly available datasets, with the only exception of the \texttt{Person} category. As \ac{ner} is the first component in the pipeline, its errors propagate to subsequent stages, amplifying their impact. This underscores the importance of domain-specific fine-tuning for \ac{ner} models---a topic further discussed in the next section (\Cref{sec:adaptation}).

Overall, these findings confirm the feasibility of extracting and semantically consolidating entity mentions from Italian court judgments, particularly in settings that implement \acl{hitl} revision of automatically generated annotations. The \ac{hitl} process used to construct the labeled dataset allowed us to estimate an average of 30 minutes to fully revise a single judgment, corresponding to approximately 320 judgments per month for one full-time annotator. This estimate provides an upper bound on the cost of human-assisted entity extraction.


\clearpage
\FloatBarrier